% Seccion "Related work" del paper de bin packing 3D (BSG_BPP).
\section{Related work}
\label{sec:related}

\pend{Statements still to be checked against the full text of the cited article: \citet{Martello2000},
\citet{Paquay2016}, \citet{Lodi2002}, \citet{Faroe2003}, \citet{Crainic2008}, \citet{Crainic2009},
\citet{Parreno2010}, \citet{Goncalves2013} (the placement procedure), \citet{Monaci2006}, the single container
loading references, and the works of the last paragraph, which are described from their abstracts.}

\paragraph{Exact methods and bounds.}
\citet{Martello2000} introduced the 3D-BPP benchmark, the lower bounds $L_1$ and $L_2$ and a
branch-and-bound method for small instances. \citet{Paquay2016} give a
mixed integer formulation for a three-dimensional bin packing problem with additional constraints.
Exact methods do not reach the sizes of the benchmark, which is why the comparisons in the
literature are made between heuristics.

\paragraph{Heuristics and metaheuristics.}
\citet{Lodi2002} proposed a tabu search framework (TS3), \citet{Faroe2003} a guided local search, \citet{Crainic2008} extreme point-based
constructive heuristics and \citet{Crainic2009} a two-level tabu search (TS2PACK).
\citet{Parreno2010} combined GRASP with variable neighbourhood descent (GVND), and
\citet{Goncalves2013} a biased random key genetic algorithm (BRKGA) with a placement procedure based
on maximal spaces; BRKGA is the reference most often used for the totals of the benchmark.
\citet{Zhu2012SD} introduced space defragmentation for packing problems, and \citet{Feng2015} two hybrid
genetic algorithms, one for small and one for large problems, that assign the items to the containers
with a largest-left-space-first strategy and pack each container with a deepest-bottom-left-with-fill
strategy. Their totals for classes~1 and 4--8 (7246 bins without rotations and 7008 with six orientations)
were obtained on ten instances per class and size that they generated with different random seeds,
with ten runs per instance (Java, a Pentium~4 at 3.4~GHz, 200 generations, about 100~s), and cannot
be compared directly with totals computed on the original instances. These methods search over
complete solutions: they move boxes between bins or reorder the sequence in which boxes are placed.
None of them recombines bins that were produced in different iterations.

\paragraph{Pattern-based methods.}
Choosing bins from a pool of patterns with a set covering or partitioning model is an established
approach. \citet{Monaci2006} applied it to the two-constraint bin packing problem and to the
two-dimensional bin packing problem without rotation, generating a small set of columns with
constructive procedures. For three-dimensional packing, \citet{Mahvash2018} use column generation
with a heuristic pricing method and a set partitioning model, and \citet{Elhedhli2019} use column
generation over layers with a set covering model and consider vertical support. In the column
generation methods the patterns come from a pricing problem built with the dual prices of the master
problem, and each pricing call is an optimisation problem in itself. Our patterns come from a search
over box priorities and from the greedy completions inside BSG, and no dual information returns to the
generator. The duals appear only after the search, to bound the pool and to prune it.

\paragraph{Single container loading.}
The loading of one bin is the single container loading problem, for which tree search methods are
the state of the art \citep{Fanslau2010,Zhu2012,Araya2014,Araya2017,Guesser2026}. BSG
\citep{Araya2014,Araya2017} is a beam search guided by a greedy algorithm, with the VCS evaluation
function; it has been extended to biobjective loading \citep{Araya2020} and to the palletization of air cargo
into several unit load devices with weight distribution constraints \citep{ArayaULD}, where the devices are
also loaded one at a time with BSG; there the objective is the volume, the devices are heterogeneous and the
total volume of the devices is minimised. \citet{Zhao2016} compare
container loading algorithms, and \citet{Bortfeldt2013} review the practical constraints, among them
vertical support.

\paragraph{Sequential and pattern-based methods for multiple containers.}
Loading the containers one at a time with an algorithm for one container is the usual approach to the multiple
container loading problem; \citet{Eley2003} notes that it is the one most often suggested in the literature, and
\citet{Martello2000} already propose a heuristic that closes one bin before opening the next. Two lines of work are
close to ours. First, \citet{LimZhang2005} and \citet{Lim2012} apply squeaky wheel optimisation
\citep{Joslin1999} to container loading. Each box type has a priority factor, initially one; in the multiple
container case the containers are filled one at a time by a layer-building heuristic with a tree search, the value of
a load is the sum over its boxes of the volume times the priority factor of the type, and after each iteration the
priority factor of every type is increased by one minus the average volume utilisation of the containers that hold
its boxes, the unloaded boxes being counted in an assumed extra container. They evaluate on the 47 instances of
Ivancic et al.\ and on 2800 instances generated from those of Bischoff and Ratcliff, not on the benchmark of
\citet{Martello2000}. Our construction and squeaky wheel loop follow the same idea, a sequential loader whose
objective is the volume weighted by a learned factor per type, with three differences: the profits start from
$v^\kappa$ with $\kappa>1$ instead of the volume, which in our experiments is the largest single effect
(Section~\ref{sec:e-ablation}); the factors are raised only for the types of the least loaded bins, multiplicatively
and with noise; and the loader is BSG. Second, \citet{Eley2003} uses a
single container algorithm to produce alternative loading patterns and selects them with a set partitioning model,
and \citet{Che2011} formulate the multiple container loading cost minimisation problem as a set cover problem whose
columns are produced by single container heuristics; \citet{Zhu2012CG} and \citet{Wei2015} generate the columns by
column generation with prototypes that approximate the pricing solutions. In our method the patterns are those that
the search builds or explores (the harvest), no pricing problem is solved, and the LP of the pool is used to stop the
recombination and to fix columns. We evaluate on the three-dimensional bin packing benchmark of \citet{Martello2000}
and against the methods designed for it.
\pend{The descriptions of \citet{Eley2003}, \citet{Che2011}, \citet{Zhu2012CG} and \citet{Wei2015} come from their
abstracts: check them in the full text, in particular which instances \citet{Che2011} and \citet{Wei2015} use
(whether they report results on the benchmark of \citet{Martello2000}), and which heuristic of \citet{Martello2000}
fills one bin at a time. Add the references to Ivancic et al.\ (1989) and Bischoff and Ratcliff (1995). State the
status of \citet{ArayaULD}.}

\paragraph{Squeaky wheel optimisation.}
Squeaky wheel optimisation \citep{Joslin1999} repeats a greedy construction and raises the priority of
the elements that the construction handled badly. \citet{Burke2011} applied it to two-dimensional strip
packing, where a penalty moves the pieces that worsened a packing to an earlier position in the next one.
In three-dimensional loading, \citet{LimZhang2005} and \citet{Lim2012} apply it to priority factors of the
box types in single and multiple container loading (see above).

\paragraph{Recent work on practical instances.}
\citet{Osaba2023} provide a benchmark dataset and an instance generator for real-world
three-dimensional bin packing. \citet{Ananno2024} propose a multi-heuristic algorithm for the
multi-container problem with real-world constraints, in which constructive heuristics build layers and
blocks and a genetic algorithm packs the remaining items. \citet{Gzara2020} extend the layer-based
column generation of \citet{Elhedhli2019} to vertical support, load bearing, planogram sequencing and weight
limits.
